% Teil II: der durchgerechnete Weg. Reihenfolge beibehalten.

\subsection*{Gegeben und gesucht}
Gegeben: $f(x)=\cos x$ auf $[0,\pi]$, mit Umkehrfunktion
$\arccos:[-1,1]\to[0,\pi]$. Gesucht: die Ableitung $\arccos'(y)$ für
$|y|<1$, und die Begründung, warum arccos bei $y=\pm1$ nicht differenzierbar
ist.

\subsection*{Ansatz}
Umkehrregel anwenden. Da $\cos$ auf $[0,\pi]$ streng monoton fallend ist und
$\cos'\neq0$ im offenen Intervall, ist arccos dort differenzierbar. Den
Zusammenhang im Einheitskreis nutzen, um $\sin$ in $\cos x$ umzuschreiben.

\subsection*{Rechnung}
Für $x=\arccos y$ gilt $\cos x=y$ und $x\in[0,\pi]$, also
$\sin x=\sqrt{1-\cos^2 x}=\sqrt{1-y^2}$ (positiv, da $x\in(0,\pi)$). Mit der
Umkehrregel:
\begin{equation}
  \arccos'(y)=\frac{1}{\cos'(\arccos y)}
  =\frac{1}{-\sin(\arccos y)}
  =-\frac{1}{\sqrt{1-y^2}}.
\end{equation}
\begin{ergebnis}
\begin{equation*}
  \arccos'(y)=-\frac{1}{\sqrt{1-y^2}},
  \qquad |y|<1.
\end{equation*}
\end{ergebnis}

\subsection*{Nichtdifferenzierbarkeit bei $y=\pm1$}
Bei $y=\pm1$ ist $x=0$ bzw.\ $x=\pi$, wo $\cos'=-\sin x=0$ ist. Der Nenner
$\sqrt{1-y^2}$ verschwindet; die Umkehrregel ist nicht anwendbar, und der
Differenzenquotient hat dort keinen endlichen Grenzwert.

\begin{kontrolle}
\textbf{Vorzeichen:} arccos fällt, daher ist die Ableitung negativ — das Minus
ist korrekt (im Gegensatz zu arcsin mit $+$).
\textbf{Grenzfall:} Für $y\to\pm1$ strebt $|\arccos'|\to\infty$
(vertikale Tangente), konsistent mit der Nichtdifferenzierbarkeit.
\textbf{Anknüpfung:} $\sin^2 x+\cos^2 x=1$ liefert die Wurzel; nur das
Vorzeichen hängt von $x\in(0,\pi)$ ab.
\end{kontrolle}

\begin{merksatz}
Die Umkehrregel $(f^{-1})'=1/(f'\circ f^{-1})$ braucht $f'\neq0$. Bei
trigonometrischen Umkehrfunktionen führt man den Sinus über
$\sin^2+\cos^2=1$ auf den Kosinus zurück; das Vorzeichen folgt aus der
Monotonie.
\end{merksatz}
